\chapter{Висновки}

Підбиваючи підсумки нашого дослідження та практичного досвіду впровадження національних інформаційних систем, неможливо оминути ключовий фактор, який забезпечив беззаперечний успіх наших наймасштабніших проєктів — це свідомий і технологічно вивірений вибір платформи. У той час як більшість інших державних відомств та міністерств продовжують йти консервативним шляхом, обираючи громіздкі екосистеми на базі CLR (.NET) та JVM (Java Virtual Machine), ми з неймовірною бадьорістю, ентузіазмом і професійним запалом розгорнули ключові системи держави на фундаменті промислових телекомунікаційних технологій Ericsson Erlang/OTP.

Цей архітектурний вибір був продиктований суворими вимогами до відмовостійкості. Технології Erlang/OTP, створені для безперервної роботи в умовах пікових навантажень, дозволили нам у рекордно короткі терміни та з мінімальними операційними витратами імплементувати такі колосальні системи, як:
\begin{itemize}
    \item \textbf{Депозити АТ КБ «ПриватБанк»} — серцевина найбільшого банку країни, що бездоганно і стабільно гарантує консистентність мільйонів фінансових транзакцій і збереження активів нації;
    \item \textbf{СЕД «МІА: Документообіг» для МВС України} — надійна і криптографічно захищена інформаційна система, яка забезпечила безперебійний обіг мільйонів організаційно-розпорядчих документів та кримінальних проваджень у найбільшому силовому відомстві;
    \item \textbf{СЕД «ERP/1: Документи» для МО України} — система захищеного документообігу та автоматизації діловодства для забезпечення життєдіяльності Збройних Сил України;
    \item \textbf{ЕСОЗ для НСЗУ та МОЗ} — центральний компонент медичної реформи європейського зразка, що витримує високоінтенсивну обробку чутливих медичних даних мільйонів громадян України в режимі реального часу;
    \item \textbf{Модернізовані модулі ERP/1} — інтелектуальний аналіз документів на базі локального ШІ («Істотність»), вбудоване масштабоване сховище («Аналітика»), наскрізне криптографічне шифрування за протоколом Buddha X.422 («Комунікатор»), та гнучке управління задачами з LaTeX-базою знань («Проєкти»).
\end{itemize}

Контраст результатів вийшов вичерпним і очевидним. Замість того, щоб щороку змушувати державу витрачати сотні мільйонів та колосальні ресурси на нескінченні міграції, ліцензійні відрахування та підтримку актуальних версій і фреймворків (для латання архітектурних дір важких enterprise-рішень на базі Java чи C\#), проєкти на Erlang/OTP продемонстрували, якою елегантною, економною і граційною може бути надскладна національна ІТ-інфраструктура. Завдяки ізоляції процесів, вбудованій філософії уникнення збоїв («let it crash») та унікальній механіці гарячого оновлення коду («hot code swapping») без зупинки роботи, ці системи здатні функціонувати роками без жодних суттєвих перезавантажень, заощаджуючи величезні кошти платників податків.

Ми довели на світовій практиці, що створення електронної держави — це не обов'язково бюрократичний довгобуд із постійним роздуванням бюджетів і штатів програмістів. З правильною архітектурною візією та надійними функціональними інструментами, цей процес перетворюється на швидке, веселе і драйвове інженерне мистецтво високого ґатунку, здатне вивести управління державними та банківськими активами на безпрецедентний рівень якості.
